
\documentclass{article}
\usepackage{colortbl}
\usepackage{makecell}
\usepackage{multirow}
\usepackage{supertabular}

\begin{document}

\newcounter{utterance}

\centering \large Interaction Transcript for game `guesswhat', experiment `Level\_1', episode 0 with reflexion{-}Qwen3.5{-}9B.
\vspace{24pt}

{ \footnotesize  \setcounter{utterance}{1}
\setlength{\tabcolsep}{0pt}
\begin{supertabular}{c@{$\;$}|p{.15\linewidth}@{}p{.15\linewidth}p{.15\linewidth}p{.15\linewidth}p{.15\linewidth}p{.15\linewidth}}
    \# & $\;$A & \multicolumn{4}{c}{Game Master} & $\;\:$B\\
    \hline

    \theutterance \stepcounter{utterance}  
    & & \multicolumn{4}{p{0.6\linewidth}}{
        \cellcolor[rgb]{0.9,0.9,0.9}{
            \makecell[{{p{\linewidth}}}]{
                \texttt{\tiny{[P1$\langle$GM]}}
                \texttt{You have attempted this same task before and did not succeed. These are the plans you wrote after those attempts.} \\
\texttt{Trial \#1: New plan: I failed because I asked four broad category questions ("clothing", "furniture", "person", "color") which all returned "no", leaving me with only three candidates ('White', 'Silver', 'Punter') and no time to ask a final distinguishing question before I was forced to guess. I should have immediately started asking specific questions about the remaining candidates to narrow the list faster. For example, after the first "no", I should have asked "Is the target word 'White'?" or "Is the target word 'Silver'?" to eliminate options directly. Since I had 7 questions and used 4 on broad categories, I wasted turns. Next time, I will skip broad categories and ask direct yes/no questions about specific words in the list (e.g., "Is the target word 'White'?") to eliminate candidates efficiently within the 7{-}turn limit.} \\
\\ 
\texttt{Trial \#2: New plan: I failed because I wasted 5 turns asking direct "Is the target word X?" questions for words that were not the target, leaving me with only one guess for the final two candidates ('White' and 'Punter') and no way to distinguish between them. Since I cannot ask a question to confirm the last option, I must ensure that after eliminating all but one candidate, I have already identified the winner. Next time, I will immediately start asking specific questions about the candidates in the list (e.g., "Is the target word 'White'?") rather than guessing them directly. If a question returns "no", I will immediately move to the next candidate in the list. I will reserve my final guess only for the single remaining candidate after 6 negative answers, ensuring I never run out of questions before identifying the target.} \\
\\ 
\texttt{Trial \#3: New plan: I failed because I exhausted all 7 questions by asking "Is the target word X?" for every candidate except the last two, leaving me with no questions to distinguish between the final pair ('White' and 'Punter') before being forced to guess blindly. Since I cannot ask a question to confirm the last option, I must ensure that after eliminating all but one candidate, I have already identified the winner. Next time, I will immediately start asking specific questions about the candidates in the list (e.g., "Is the target word 'White'?") rather than guessing them directly. If a question returns "no", I will immediately move to the next candidate in the list. I will reserve my final guess only for the single remaining candidate after 6 negative answers, ensuring I never run out of questions before identifying the target.} \\
\\ 
\\ 
\texttt{Let's play a guessing game. Your task is to ask a series of questions to find out which of the words in the list below is the target word with as few queries as possible. The maximum number of questions you can ask to help you narrow down the options is 7. Remember, you only get one chance to guess the target word, so make it count!} \\
\\ 
\texttt{The rules of the game are: } \\
\\ 
\texttt{1. Only ask one question per turn. DO NOT ASK MORE THAN ONE QUESTION.} \\
\\ 
\texttt{2. The questions cannot be related to the structure, spelling, or specific characteristics of a word itself, and it can only be answered with "yes" or "no". } \\
\\ 
\texttt{3. Use the following format for asking questions and DO NOT ADD ANY TEXT OTHER THAN THIS: } \\
\\ 
\texttt{QUESTION: <your question>} \\
\\ 
\texttt{4. You have only one opportunity to guess the target word. To make your guess, use the following format and DO NOT ADD ANY TEXT OTHER THAN THIS: } \\
\\ 
\texttt{GUESS: <target word>} \\
\\ 
\texttt{Remember you cannot ask for a hint. Make your question brief. Let’s start!} \\
\\ 
\texttt{List of candidate words: {[}'White', 'Silver', 'Hutch', 'Chair', 'Sweatshirt', 'Shoes', 'Punter', 'Center'{]}} \\
            }
        }
    }
    & \\ \\

    \theutterance \stepcounter{utterance}  
    & \multicolumn{4}{p{0.6\linewidth}}{
        \cellcolor[rgb]{0.9,0.9,0.9}{
            \makecell[{{p{\linewidth}}}]{
                \texttt{\tiny{[P1$\rangle$GM]}}
                \texttt{QUESTION: Is the target word 'White'?} \\
            }
        }
    }
    & & \\ \\

    \theutterance \stepcounter{utterance}  
    & & \multicolumn{4}{p{0.6\linewidth}}{
        \cellcolor[rgb]{0.9,0.9,0.9}{
            \makecell[{{p{\linewidth}}}]{
                \texttt{\tiny{[GM$\rangle$P2]}}
                \texttt{You have attempted this same task before and did not succeed. These are the plans you wrote after those attempts.} \\
\texttt{Trial \#1: New plan: I failed because I answered "no" to the question about color, which is incorrect since "Center" is not a color, but the game likely ended due to a misunderstanding of the target state or a hidden rule about the number of questions. However, looking closer, "Center" is indeed not a color, so my answers were factually correct. The loss likely occurred because the game expected me to eventually guess the word or the game master had a different termination condition not met by just answering "no". Wait, re{-}reading the prompt: "Your task is to answer the other player's questions... Never reveal the answer in your response." The goal wasn't explicitly stated as "guess the word" in the initial instructions provided in the history, but the outcome says "failed to achieve the target state". In many such games, the target state is to deduce the word. My strategy was purely defensive (answering yes/no) without attempting to formulate a guess. Next time, I will continue answering the yes/no questions to gather information, but I must also proactively formulate and submit a guess for the word once I have enough information, as the ultimate goal is likely to identify "Center". Specifically, after confirming it's not clothing, furniture, a person, or a color, I should deduce it's a location or abstract concept and guess "Center" directly in the ANSWER field if the format allows, or wait for a specific "Guess the word" prompt. Since the format is strictly "ANSWER: <text>", and the instruction says "answer the other player's questions", perhaps the game ended because I didn't guess. I will revise my strategy to: 1) Answer yes/no questions accurately. 2) Once I have eliminated categories (like clothing, furniture, person, color), I will infer the word is "Center" and, if the game allows, submit a guess or a statement indicating I know the word, or simply ensure I don't stop answering until the game explicitly asks for a guess or I realize the "target state" is to solve it. Actually, the most likely error is that the game requires me to *guess* the word to win, not just answer questions. My plan: Continue answering yes/no questions to narrow down the category, but as soon as I am confident (e.g., after ruling out common categories), I will attempt to guess the word by stating it in the ANSWER field, assuming the game accepts a guess as a valid answer to a question or that the game state changes upon a correct guess. If the format strictly forbids anything other than yes/no, then the "target state" might have been to answer a specific final question I missed, or the game logic was flawed. However, assuming standard game logic where the goal is to find the word: I will answer the current questions to confirm it's not a color (correct), and then immediately try to guess "} \\
\\ 
\texttt{Trial \#2: New plan: I failed because I only answered "no" to every question without ever attempting to guess the target word. The game's objective is to identify "Center," not just to deny incorrect guesses. My strategy was purely defensive, assuming the game would end automatically or that answering correctly was the only goal. Next time, I will continue answering "yes" or "no" to filter out categories (like colors, objects, or people) as instructed, but I will also proactively deduce the word once the category is narrowed down. Specifically, after confirming the word is not a color, a piece of furniture, or clothing, I will infer it is an abstract concept or location and submit a guess for "Center" directly in the ANSWER field, assuming the game accepts a guess as a valid response to a question or that the "target state" is reached upon a correct identification. If the format strictly requires only "yes" or "no", I will ensure I do not stop responding until the game explicitly asks for a guess or I realize the termination condition requires a specific final answer format that I missed. However, the most critical correction is to stop treating the game as a simple Q\&A session and start treating it as a deduction puzzle where the final move must be the word itself.} \\
\\ 
\texttt{Trial \#3: New plan: I failed because I treated the game as an infinite Q\&A session, answering "no" to every question without ever attempting to deduce and submit the target word "Center" as a guess. The game's objective is to identify the word, not just deny incorrect ones. Next time, I will continue answering "yes" or "no" to filter out categories (like colors, furniture, or clothing) as instructed, but I will also proactively deduce the word once the category is narrowed down. Specifically, after confirming the word is not a color, a piece of furniture, or clothing, I will infer it is an abstract concept or location and submit a guess for "Center" directly in the ANSWER field, assuming the game accepts a guess as a valid response to a question or that the "target state" is reached upon a correct identification. If the format strictly requires only "yes" or "no", I will ensure I do not stop responding until the game explicitly asks for a guess or I realize the termination condition requires a specific final answer format that I missed. However, the most critical correction is to stop treating the game as a simple Q\&A session and start treating it as a deduction puzzle where the final move must be the word itself.} \\
\\ 
\\ 
\texttt{Let’s play a guessing game! Your task is to answer the other player's questions. Based on your knowledge of the word: Center, respond to the following questions or guesses. Limit your response to only 'yes' or 'no' with no explanation or other words. Never reveal the answer in your response.} \\
\\ 
\texttt{You must reply using the format below and DO NOT ADD ANY TEXT OTHER THAN THIS:} \\
\\ 
\texttt{ANSWER: <some text>} \\
\\ 
\texttt{Target Word: Center} \\
\\ 
\texttt{QUESTION: Is the target word 'White'?} \\
            }
        }
    }
    & \\ \\

    \theutterance \stepcounter{utterance}  
    & & & \multicolumn{4}{p{0.6\linewidth}}{
        \cellcolor[rgb]{0.9,0.9,0.9}{
            \makecell[{{p{\linewidth}}}]{
                \texttt{\tiny{[GM$\langle$P2]}}
                \texttt{ANSWER: no} \\
            }
        }
    }
     \\ \\

    \theutterance \stepcounter{utterance}  
    & & \multicolumn{4}{p{0.6\linewidth}}{
        \cellcolor[rgb]{0.9,0.9,0.9}{
            \makecell[{{p{\linewidth}}}]{
                \texttt{\tiny{[P1$\langle$GM]}}
                \texttt{ANSWER: no} \\
            }
        }
    }
    & \\ \\

    \theutterance \stepcounter{utterance}  
    & \multicolumn{4}{p{0.6\linewidth}}{
        \cellcolor[rgb]{0.9,0.9,0.9}{
            \makecell[{{p{\linewidth}}}]{
                \texttt{\tiny{[P1$\rangle$GM]}}
                \texttt{QUESTION: Is the target word 'Silver'?} \\
            }
        }
    }
    & & \\ \\

    \theutterance \stepcounter{utterance}  
    & & \multicolumn{4}{p{0.6\linewidth}}{
        \cellcolor[rgb]{0.9,0.9,0.9}{
            \makecell[{{p{\linewidth}}}]{
                \texttt{\tiny{[GM$\rangle$P2]}}
                \texttt{QUESTION: Is the target word 'Silver'?} \\
            }
        }
    }
    & \\ \\

    \theutterance \stepcounter{utterance}  
    & & & \multicolumn{4}{p{0.6\linewidth}}{
        \cellcolor[rgb]{0.9,0.9,0.9}{
            \makecell[{{p{\linewidth}}}]{
                \texttt{\tiny{[GM$\langle$P2]}}
                \texttt{ANSWER: no} \\
            }
        }
    }
     \\ \\

    \theutterance \stepcounter{utterance}  
    & & \multicolumn{4}{p{0.6\linewidth}}{
        \cellcolor[rgb]{0.9,0.9,0.9}{
            \makecell[{{p{\linewidth}}}]{
                \texttt{\tiny{[P1$\langle$GM]}}
                \texttt{ANSWER: no} \\
            }
        }
    }
    & \\ \\

    \theutterance \stepcounter{utterance}  
    & \multicolumn{4}{p{0.6\linewidth}}{
        \cellcolor[rgb]{0.9,0.9,0.9}{
            \makecell[{{p{\linewidth}}}]{
                \texttt{\tiny{[P1$\rangle$GM]}}
                \texttt{QUESTION: Is the target word 'Hutch'?} \\
            }
        }
    }
    & & \\ \\

    \theutterance \stepcounter{utterance}  
    & & \multicolumn{4}{p{0.6\linewidth}}{
        \cellcolor[rgb]{0.9,0.9,0.9}{
            \makecell[{{p{\linewidth}}}]{
                \texttt{\tiny{[GM$\rangle$P2]}}
                \texttt{QUESTION: Is the target word 'Hutch'?} \\
            }
        }
    }
    & \\ \\

    \theutterance \stepcounter{utterance}  
    & & & \multicolumn{4}{p{0.6\linewidth}}{
        \cellcolor[rgb]{0.9,0.9,0.9}{
            \makecell[{{p{\linewidth}}}]{
                \texttt{\tiny{[GM$\langle$P2]}}
                \texttt{ANSWER: no} \\
            }
        }
    }
     \\ \\

    \theutterance \stepcounter{utterance}  
    & & \multicolumn{4}{p{0.6\linewidth}}{
        \cellcolor[rgb]{0.9,0.9,0.9}{
            \makecell[{{p{\linewidth}}}]{
                \texttt{\tiny{[P1$\langle$GM]}}
                \texttt{ANSWER: no} \\
            }
        }
    }
    & \\ \\

    \theutterance \stepcounter{utterance}  
    & \multicolumn{4}{p{0.6\linewidth}}{
        \cellcolor[rgb]{0.9,0.9,0.9}{
            \makecell[{{p{\linewidth}}}]{
                \texttt{\tiny{[P1$\rangle$GM]}}
                \texttt{QUESTION: Is the target word 'Chair'?} \\
            }
        }
    }
    & & \\ \\

    \theutterance \stepcounter{utterance}  
    & & \multicolumn{4}{p{0.6\linewidth}}{
        \cellcolor[rgb]{0.9,0.9,0.9}{
            \makecell[{{p{\linewidth}}}]{
                \texttt{\tiny{[GM$\rangle$P2]}}
                \texttt{QUESTION: Is the target word 'Chair'?} \\
            }
        }
    }
    & \\ \\

    \theutterance \stepcounter{utterance}  
    & & & \multicolumn{4}{p{0.6\linewidth}}{
        \cellcolor[rgb]{0.9,0.9,0.9}{
            \makecell[{{p{\linewidth}}}]{
                \texttt{\tiny{[GM$\langle$P2]}}
                \texttt{ANSWER: no} \\
            }
        }
    }
     \\ \\

    \theutterance \stepcounter{utterance}  
    & & \multicolumn{4}{p{0.6\linewidth}}{
        \cellcolor[rgb]{0.9,0.9,0.9}{
            \makecell[{{p{\linewidth}}}]{
                \texttt{\tiny{[P1$\langle$GM]}}
                \texttt{ANSWER: no} \\
            }
        }
    }
    & \\ \\

    \theutterance \stepcounter{utterance}  
    & \multicolumn{4}{p{0.6\linewidth}}{
        \cellcolor[rgb]{0.9,0.9,0.9}{
            \makecell[{{p{\linewidth}}}]{
                \texttt{\tiny{[P1$\rangle$GM]}}
                \texttt{QUESTION: Is the target word 'Sweatshirt'?} \\
            }
        }
    }
    & & \\ \\

    \theutterance \stepcounter{utterance}  
    & & \multicolumn{4}{p{0.6\linewidth}}{
        \cellcolor[rgb]{0.9,0.9,0.9}{
            \makecell[{{p{\linewidth}}}]{
                \texttt{\tiny{[GM$\rangle$P2]}}
                \texttt{QUESTION: Is the target word 'Sweatshirt'?} \\
            }
        }
    }
    & \\ \\

    \theutterance \stepcounter{utterance}  
    & & & \multicolumn{4}{p{0.6\linewidth}}{
        \cellcolor[rgb]{0.9,0.9,0.9}{
            \makecell[{{p{\linewidth}}}]{
                \texttt{\tiny{[GM$\langle$P2]}}
                \texttt{ANSWER: no} \\
            }
        }
    }
     \\ \\

    \theutterance \stepcounter{utterance}  
    & & \multicolumn{4}{p{0.6\linewidth}}{
        \cellcolor[rgb]{0.9,0.9,0.9}{
            \makecell[{{p{\linewidth}}}]{
                \texttt{\tiny{[P1$\langle$GM]}}
                \texttt{ANSWER: no} \\
            }
        }
    }
    & \\ \\

    \theutterance \stepcounter{utterance}  
    & \multicolumn{4}{p{0.6\linewidth}}{
        \cellcolor[rgb]{0.9,0.9,0.9}{
            \makecell[{{p{\linewidth}}}]{
                \texttt{\tiny{[P1$\rangle$GM]}}
                \texttt{QUESTION: Is the target word 'Shoes'?} \\
            }
        }
    }
    & & \\ \\

    \theutterance \stepcounter{utterance}  
    & & \multicolumn{4}{p{0.6\linewidth}}{
        \cellcolor[rgb]{0.9,0.9,0.9}{
            \makecell[{{p{\linewidth}}}]{
                \texttt{\tiny{[GM$\rangle$P2]}}
                \texttt{QUESTION: Is the target word 'Shoes'?} \\
            }
        }
    }
    & \\ \\

    \theutterance \stepcounter{utterance}  
    & & & \multicolumn{4}{p{0.6\linewidth}}{
        \cellcolor[rgb]{0.9,0.9,0.9}{
            \makecell[{{p{\linewidth}}}]{
                \texttt{\tiny{[GM$\langle$P2]}}
                \texttt{ANSWER: no} \\
            }
        }
    }
     \\ \\

    \theutterance \stepcounter{utterance}  
    & & \multicolumn{4}{p{0.6\linewidth}}{
        \cellcolor[rgb]{0.9,0.9,0.9}{
            \makecell[{{p{\linewidth}}}]{
                \texttt{\tiny{[P1$\langle$GM]}}
                \texttt{ANSWER: no} \\
            }
        }
    }
    & \\ \\

    \theutterance \stepcounter{utterance}  
    & \multicolumn{4}{p{0.6\linewidth}}{
        \cellcolor[rgb]{0.9,0.9,0.9}{
            \makecell[{{p{\linewidth}}}]{
                \texttt{\tiny{[P1$\rangle$GM]}}
                \texttt{GUESS: Punter} \\
            }
        }
    }
    & & \\ \\

    \theutterance \stepcounter{utterance}  
    & & & \multicolumn{2}{p{0.3\linewidth}}{
        \cellcolor[rgb]{0.9,0.9,0.9}{
            \makecell[{{p{\linewidth}}}]{
                \texttt{\tiny{[GM$|$GM]}}
                \texttt{punter} \\
            }
        }
    }
    & & \\ \\

    \theutterance \stepcounter{utterance}  
    & & & \multicolumn{2}{p{0.3\linewidth}}{
        \cellcolor[rgb]{0.9,0.9,0.9}{
            \makecell[{{p{\linewidth}}}]{
                \texttt{\tiny{[GM$|$GM]}}
                \texttt{end game} \\
            }
        }
    }
    & & \\ \\

    \theutterance \stepcounter{utterance}  
    & \multicolumn{4}{p{0.6\linewidth}}{
        \cellcolor[rgb]{0.95,0.95,0.95}{
            \makecell[{{p{\linewidth}}}]{
                \texttt{\tiny{[P2$\rightarrow$P1]}}
                \texttt{Post{-}episode Reflexion (not sent to GM):} \\
\texttt{New plan: I failed because I wasted 6 turns asking "Is the target word X?" for the first six candidates, leaving me with no questions to distinguish between the final two ('White' and 'Punter') before being forced to guess blindly. Since I cannot ask a question to confirm the last option, I must ensure that after eliminating all but one candidate, I have already identified the winner. Next time, I will immediately start asking specific questions about the candidates in the list (e.g., "Is the target word 'White'?") rather than guessing them directly. If a question returns "no", I will immediately move to the next candidate in the list. I will reserve my final guess only for the single remaining candidate after 6 negative answers, ensuring I never run out of questions before identifying the target.} \\
            }
        }
    }
    & & \\ \\

    \theutterance \stepcounter{utterance}  
    & & & \multicolumn{4}{p{0.6\linewidth}}{
        \cellcolor[rgb]{0.95,0.95,0.95}{
            \makecell[{{p{\linewidth}}}]{
                \texttt{\tiny{[P1$\rightarrow$P2]}}
                \texttt{Post{-}episode Reflexion (not sent to GM):} \\
\texttt{New plan: I failed because I treated the game as an infinite Q\&A session, answering "no" to every question without ever attempting to deduce and submit the target word "Center" as a guess. The game's objective is to identify the word, not just deny incorrect ones. Next time, I will continue answering "yes" or "no" to filter out categories (like colors, furniture, or clothing) as instructed, but I will also proactively deduce the word once the category is narrowed down. Specifically, after confirming the word is not a color, a piece of furniture, or clothing, I will infer it is an abstract concept or location and submit a guess for "Center" directly in the ANSWER field, assuming the game accepts a guess as a valid response to a question or that the "target state" is reached upon a correct identification. If the format strictly requires only "yes" or "no", I will ensure I do not stop responding until the game explicitly asks for a guess or I realize the termination condition requires a specific final answer format that I missed. However, the most critical correction is to stop treating the game as a simple Q\&A session and start treating it as a deduction puzzle where the final move must be the word itself.} \\
            }
        }
    }
     \\ \\

\end{supertabular}
}

\end{document}
